% संपूर्ण मराठी ग्रंथातील प्रकरण 36: विन्यासमीमांसा
% हा संपादनयोग्य स्रोतखंड आहे. संकलनासाठी संपूर्ण स्रोतसंग्रहातील
% अक्षररूपे, पूर्वभाग, चिन्हव्याख्या आणि आकृती-स्रोत आवश्यक आहेत.
% मूळ: Open Logic Project; मजकूर आणि रूपांतर CC BY 4.0.
% यंत्रानुवाद, दुरुस्ती आणि तपासणी: OpenAI Codex —
% GPT-5.6 Sol आणि GPT-6 Sol, दोन्ही Ultra effort.
% दुरुस्ती-पुनरावलोकन: GPT-6.1 Sol, Ultra effort.
\chapter{विन्यासमीमांसा}\label{lam:syn:chap}








\section{पदे}\label{lam:syn:trm:sec}

लॅम्डा कलनातील पदे $\Obj{v_0}$, $\Obj{v_1}$, \dots या
अमर्याद चरांपासून, ~“$\lambd$” या चिन्हापासून आणि कंसांपासून
विगमनाधारित रीतीने घडवली जातात. चरे दर्शवण्यासाठी आपण
$x$, $y$, $z$, \dots{} आणि पदे दर्शवण्यासाठी
$M$, $N$, $P$, \dots{} वापरू.

\begin{defn}[पदे] \label{lam:syn:trm:defn:term}
लॅम्डा कलनातील \emph{पदांचा} संच पुढीलप्रमाणे
विगमनाधारित रीतीने परिभाषित करतात:
\begin{enumerate}
  \item \label{lam:syn:trm:defn:term-var} $x$ हे चर असेल, तर $x$ हे
    पद आहे.
  \item \label{lam:syn:trm:defn:term-abs} $x$ हे चर आणि $M$ हे
    पद असेल, तर $(\lambd[x][M])$ हे पद आहे.
  \item \label{lam:syn:trm:defn:term-app} $M$ आणि $N$ ही दोन्ही
    पदे असतील, तर $(MN)$ हे पद आहे.
\end{enumerate}
\end{defn}

पद $(\lambd[x][M])$ हे \ref{lam:syn:trm:defn:term-abs} नुसार
घडवले असेल, तर ते \emph{अमूर्तीकरणाचे} फलित आहे
असे म्हणतो; आणि $\lambd[x]$ मधील $x$ ला
\emph{प्राचल} म्हणतात. \ref{lam:syn:trm:defn:term-app}
नुसार घडवलेले पद $(MN)$ हे \emph{उपयोजनाचे}
फलित असते.

वर परिभाषित पदांमध्ये सर्व कंस स्पष्टपणे लिहिलेले
असतात. उदाहरणार्थ,
$(\lambd[x][((\lambd[x][x])(\lambd[x][(xx)]))])$
या पदावरून दिसते की हे खूपच अवजड होऊ शकते. कंस
टाळण्यासाठी आपण काही रूढी मांडू. मात्र अधिकृत
व्याख्येमुळे \ref{lam:syn:trm:defn:term} नुसार पद कसे
घडले हे ठरवणे सोपे जाते. उदाहरणार्थ,
$(\lambd[x][((\lambd[x][x])(\lambd[x][(xx)]))])$
हे पद घडवण्याची शेवटची पायरी अमूर्तीकरणच असली
पाहिजे; त्यात प्राचल हे~$x$ आहे.
ते $((\lambd[x][x])(\lambd[x][(xx)]))$ या पदावर
अमूर्तीकरण करून मिळते; हे पद दोन पदांच्या
उपयोजनातून घडले आहे. त्या दोन पदांपैकी प्रत्येक
अमूर्तीकरणातून घडले आहे; आणि पुढे असेच चालते.

\begin{prob}
$(\lambd[g][(\lambd[x][(g (x x))])
  (\lambd[x][(g (x x))])])$ हे पद कसे घडले ते वर्णन करा.
\end{prob}











\section{एकमेव वाचनीयता}\label{lam:syn:unq:sec}

प्रत्येक पद घडवण्याचा एकमेव मार्ग असतो का, असा प्रश्न
पडू शकतो; आणि तसेच आहे. प्रत्येक लॅम्डा पद घडवण्याची
आणि त्याचा अर्थ लावण्याची एकच पद्धत आहे. पुढील चर्चेत
\emph{रचनाप्रक्रिया} म्हणजे \ref{lam:syn:trm:defn:term}
मधील रचनानियम एकदा किंवा अनेकदा वापरून पद घडवण्याची
प्रक्रिया.

\begin{lem}\label{lam:syn:unq:lem:term-start}
पदाची सुरुवात एकतर चराने किंवा कंसाने होते.
\end{lem}

\begin{proof}
\ref{lam:syn:trm:defn:term} नुसार घडवले असेल तरच एखादी
गोष्ट पद ठरते. ते \ref{lam:syn:trm:defn:term-var} चे फलित
असेल, तर ते चर असलेच पाहिजे. ते
\ref{lam:syn:trm:defn:term-abs} किंवा
\ref{lam:syn:trm:defn:term-app} चे फलित असेल, तर
त्याची सुरुवात कंसाने होते.
\end{proof}

\begin{lem}\label{lam:syn:unq:lem:app-start}
उपयोजनाच्या फलिताची सुरुवात एकतर दोन कंसांनी किंवा
एक कंस आणि त्यानंतर चर याने होते.
\end{lem}

\begin{proof}
$M$ हे उपयोजनाचे फलित असेल, तर ते $(PQ)$ या रूपात
असते; म्हणून त्याची सुरुवात कंसाने होते. $P$ हे पद
असल्यामुळे \ref{lam:syn:unq:lem:term-start} नुसार त्याची
सुरुवात कंसाने किंवा चराने होते.
\end{proof}

\begin{lem}\label{lam:syn:unq:lem:initial}
पदाचा कोणताही उचित आरंभीचा भाग स्वतः पद नसतो.
\end{lem}

\begin{prob}
पदांच्या लांबीवरील विगमनाने
\ref{lam:syn:unq:lem:initial} सिद्ध करा.
\end{prob}

\begin{prop}[एकमेव वाचनीयता] \label{lam:syn:unq:prop:unq}
प्रत्येक पदाची रचनाप्रक्रिया एकमेव असते. दुसऱ्या
शब्दांत, एखाद्या रचनाप्रक्रियेने पद $M$ घडत असेल,
तर तेच पद घडवू शकणारी ती एकमेव रचनाप्रक्रिया असते.
\end{prop}

\begin{proof}
पदांच्या रचनाप्रक्रियेवरील विगमनाने हे सिद्ध करू.
  \begin{enumerate}
    \item $M$ हे $x$ या रूपात आहे, जिथे $x$ हे
      एखादे चर आहे. अमूर्तीकरण आणि उपयोजन यांची
      फलिते नेहमी कंसाने सुरू होतात, म्हणून $M$
      घडवताना त्यांचा वापर झालेला असू शकत नाही.
      त्यामुळे $M$ ची रचनाप्रक्रिया
      \ref{lam:syn:trm:defn:term}\ref{lam:syn:trm:defn:term-var}
      ची एकच पायरी असली पाहिजे.
    \item $M$ हे $(\lambd[x][N])$ या रूपात आहे,
      जिथे $x$ हे चर आणि $N$ हे पद आहे. ते
      \ref{lam:syn:trm:defn:term}\ref{lam:syn:trm:defn:term-var}
      नुसार घडलेले असू शकत नाही, कारण ते एकटे
      चर नाही. \ref{lam:syn:unq:lem:app-start} नुसार ते
      उपयोजनाचेही फलित नाही. त्यामुळे $M$ हे
      फक्त~$N$ वरील अमूर्तीकरणाचे फलित असू शकते.
      विगमन गृहीतकानुसार $N$ ची रचनाप्रक्रियाही
      एकमेव आहे.
    \item $M$ हे $(PQ)$ या रूपात आहे, जिथे $P$
      आणि $Q$ ही पदे आहेत. त्याची सुरुवात कंसाने
      होत असल्याने ते
      \ref{lam:syn:trm:defn:term}\ref{lam:syn:trm:defn:term-var}
      नुसारही घडलेले असू शकत नाही.
      \ref{lam:syn:unq:lem:term-start} नुसार $P$ ची सुरुवात
      $\lambd$ ने होऊ शकत नाही; म्हणून $(PQ)$ हे
      अमूर्तीकरणाचे फलितही असू शकत नाही. आता $M$
      उपयोजनाने घडवण्याचा आणखी एक मार्ग आहे असे
      समजा; उदाहरणार्थ, ते $(P'Q')$ या रूपातही
      आहे. सुरुवातीची पदे वेगळी असतील, तर $P$
      हा $P'$ चा उचित आरंभीचा भाग असेल किंवा
      उलट असेल; \ref{lam:syn:unq:lem:initial} नुसार हे
      अशक्य आहे. ती एकच असतील, तर उरलेले भागही
      एकच असतात. म्हणून $P$ आणि $Q$ एकमेव
      रीतीने ठरतात आणि विगमन गृहीतकानुसार $P$
      आणि $Q$ यांच्या रचनाप्रक्रियाही एकमेव असतात.
  \end{enumerate}
\end{proof}

वरील प्रतिज्ञेचा अधिक वाचनीय अर्थ असा आहे:
\begin{prop}
  पद $M$ पुढीलपैकी एकाच रूपात असू शकते:
  \begin{enumerate}
    \item $x$, जिथे $x$ हे $M$ ने एकमेव रीतीने
      ठरणारे चर आहे.
    \item $(\lambd[x][N])$, जिथे $x$ हे चर आणि
      $N$ हे दुसरे पद आहे; दोन्ही $M$ ने एकमेव
      रीतीने ठरतात.
    \item $(PQ)$, जिथे $P$ आणि $Q$ ही दोन पदे
      $M$ ने एकमेव रीतीने ठरतात.
  \end{enumerate}
\end{prop}











\section{संक्षिप्त विन्यास}\label{lam:syn:abb:sec}

\ref{lam:syn:trm:defn:term} मध्ये परिभाषित पदे लिहिताना
कधीकधी अवजड वाटतात; म्हणून अधिक संक्षिप्त विन्यास
मांडणे उपयुक्त आहे. अर्थात, या संक्षिप्त संकेतनातील
पदेही एकमेव रीतीने वाचता येतील याची काळजी
घ्यावी लागेल. \emph{संक्षिप्त पदे} म्हणून ओळखली
जाणारी एक प्रचलित पद्धत पुढीलप्रमाणे आहे.

\begin{enumerate}
\item कंस वगळले असता उपयोजन डावीकडून उजवीकडे
  होते. उदाहरणार्थ, $M$, $N$, $P$ आणि~$Q$
  ही पदे असतील, तर $MNPQ$ हे $(((MN)P)Q)$
  चे संक्षिप्त रूप आहे.
\item पुन्हा, कंस वगळले असता लॅम्डा अमूर्तीकरणाला
  शक्य तितकी व्यापक व्याप्ती द्यायची. उदाहरणार्थ,
  $\lambd[x][MNP]$ याचा अर्थ
  $(\lambd[x][((MN)P)])$ असा घ्यायचा.
\item एका लॅम्डाने अनेक चरे अमूर्त करता येतात.
  उदाहरणार्थ, $\lambd[xyz][M]$ हे
  $\lambd[x][\lambd[y][\lambd[z][M]]]$
  चे संक्षिप्त रूप आहे.
\end{enumerate}

उदाहरणार्थ,
\[\lambd[xy][xxyx \lambd[z][xz]]
\]
हे पुढीलचे संक्षिप्त रूप आहे:
\[(\lambd[x][(\lambd[y][((((xx)y)x)(\lambd[z][(xz)]))])]).
\]

\begin{prob}
संक्षिप्त पद
$\lambd[g][(\lambd[x][g (x x)])
  \lambd[x][g (x x)]]$
विस्तारून लिहा.
\end{prob}











\section{मुक्त चरे}\label{lam:syn:fv:sec}

लॅम्डा कलनात फलनांचा अभ्यास केला जातो, आणि लॅम्डा अमूर्तीकरणातून
फलने मिळतात. सहज समजुतीने, $\lambd[x][M]$ हे असे फलन आहे की
फलनाचा फलसाधक~$x$ ला दिल्यावर त्याची मूल्ये~$M$ ने
ठरतात. पण~$M$ मधील $x$ ची प्रत्येक आस्थिती यासाठी संबंधित नसते:
$M$ मध्ये आणखी एखादे अमूर्तीकरण $\lambd[x][N]$ असेल, तर~$N$
मधील $x$ च्या आस्थिती $\lambd[x][N]$ शी संबंधित असतात;
$\lambd[x][M]$ शी नाही. म्हणून $\lambd[x][M]$ मधील
लॅम्डा अमूर्तीकरण $\lambd[x]$ हे~$M$ मधील $x$ च्या अशा
आस्थिती \emph{बद्ध करते} ज्या दुसऱ्या लॅम्डा अमूर्तीकरणाने
आधीच बद्ध केलेल्या नाहीत---म्हणजे~$M$ मधील $x$ च्या
\emph{मुक्त} आस्थिती.

\begin{defn}[व्याप्ती]
पद~$N$ मध्ये $\lambd[x][M]$ आढळत असेल, तर त्यातील~$M$
ची संबंधित आस्थिती ही~$\lambd[x]$ ची
\emph{व्याप्ती} असते.
\end{defn}


\begin{defn}[मुक्त आणि बद्ध आस्थिती]
  पद~$M$ मधील चर $x$ ची एखादी आस्थिती $\lambd[x]$ च्या
  व्याप्तीत नसेल तर ती \emph{मुक्त} असते; अन्यथा ती
  \emph{बद्ध} असते. $\lambd[x][M]$ मधील चर~$x$ ची
  एखादी आस्थिती सुरुवातीच्या $\lambd[x]$ ने बद्ध केलेली
  असते, जर आणि फक्त जर~$M$ मधील $x$ ची ती आस्थिती मुक्त असते.
\end{defn}

\begin{ex}
  $\lambd[x][x y]$ मध्ये $x$ आणि $y$ हे दोन्ही $\lambd[x]$
  च्या व्याप्तीत आहेत; म्हणून $x$ हे $\lambd[x]$ ने बद्ध होते.
  $y$ कोणत्याही $\lambd[y]$ च्या व्याप्तीत नसल्याने ते मुक्त
  असते. $\lambd[x][x x]$ मध्ये $x$ च्या दोन्ही आस्थिती
  $\lambd[x]$ ने बद्ध होतात, कारण $xx$ मध्ये त्या दोन्ही मुक्त
  आहेत. $((\lambd[x][xx])x)$ मध्ये~$x$ ची शेवटची आस्थिती
  मुक्त आहे, कारण ती कोणत्याही $\lambd[x]$ च्या व्याप्तीत नाही.
  $\lambd[x][(\lambd[x][x])x]$ मध्ये पहिल्या $\lambd[x]$
  ची व्याप्ती $(\lambd[x][x])x$ आहे आणि दुसऱ्या $\lambd[x]$
  ची व्याप्ती शेवटून दुसरी असलेली~$x$ ची आस्थिती आहे.
  $(\lambd[x][x])x$ मध्ये $x$ ची शेवटची आस्थिती मुक्त,
  तर शेवटून दुसरी बद्ध आहे. म्हणून
  $\lambd[x][(\lambd[x][x])x]$ मधील $x$ ची शेवटून दुसरी
  आस्थिती दुसऱ्या $\lambd[x]$ ने, आणि शेवटची आस्थिती
  पहिल्या $\lambd[x]$ ने बद्ध होते.
\end{ex}

पद $P$ मधील चरांच्या सर्व आस्थिती तपासून मुक्त चरांचा संच
मिळवता येतो. हा संच $\FV{P}$ असा दर्शवतात; त्याची
स्वाभाविक व्याख्या पुढीलप्रमाणे आहे:

\begin{defn}[पदाचे मुक्त चरे] \label{lam:syn:fv:def:fv}
  पदाच्या \emph{मुक्त चरांचा} संच विगमनाने पुढीलप्रमाणे परिभाषित करतात:
  \begin{enumerate}
    \item $\FV{x} = \{x\}$ \label{lam:syn:fv:def:fv1}
    \item $\FV{\lambd[x][N]} = \FV{N} \setminus \{x\}$ \label{lam:syn:fv:def:fv2}
    \item $\FV{PQ} = \FV{P} \cup \FV{Q}$ \label{lam:syn:fv:def:fv3}
  \end{enumerate}
\end{defn}

\begin{prob}
  \begin{enumerate}
  \item या पदातील $\lambd[g]$ आणि दोन $\lambd[x]$ यांच्या
    व्याप्ती ओळखा: $\lambd[g][(\lambd[x][g (x x)]) \lambd[x][g (x x)]]$.
  \item $\lambd[g][(\lambd[x][g (x x)]) \lambd[x][g (x x)]]$
    मधील चरांच्या सर्व आस्थिती बद्ध आहेत का? असल्यास,
    त्यांना अनुक्रमे कोणती अमूर्तीकरणे बद्ध करतात?
    \item $\FV{\lambd[x][(\lambd[y][(\lambd[z][x y]) z]) y]}$
      काढा.
  \end{enumerate}
\end{prob}

\begin{explain}
मुक्त चर म्हणजे जणू बाह्य जगाचा (\emph{परिसराचा}) संदर्भ.
मुक्त चरे असलेले पद अंशतःच निश्चित झालेले पद मानता येते,
कारण त्याचे वर्तन आपण परिसर कसा ठरवतो यावर अवलंबून असते.
उदाहरणार्थ, $\lambd[x][f x]$ हे पद फलसाधक~$x$ स्वीकारून
त्यावर $f$ लावून मिळणारा परिणाम देते; येथे चर~$f$ मुक्त आहे.
या पदाचे मूल्य ते ज्या परिसरात आहे त्यावर, विशेषतः त्या
परिसरातील $f$ च्या मूल्यावर, अवलंबून असते.

या पदावर अमूर्तीकरण केल्यास $\lambd[f][\lambd[x][f
    x]]$ मिळते. आता हे पद परिसरातील चर $f$ वर अवलंबून
नसते, कारण ते दोन फलसाधक स्वीकारून पहिला दुसऱ्यावर
लावल्याचा परिणाम देणारे फलन दर्शवते. परिसरातील $f$
बदलला तरी या पदाच्या वर्तनावर परिणाम होत नाही: हे पद
फक्त फलसाधक म्हणून दिलेले मूल्य वापरते, परिसरातील
$f$ चे मूल्य वापरत नाही.
\end{explain}

\begin{defn}[बंद पद, संयोजक]
  मुक्त चरे नसलेल्या पदाला \emph{बंद पद} किंवा
  \emph{संयोजक} म्हणतात.
\end{defn}

\begin{lem}\label{lam:syn:fv:lem:fv}
  \begin{enumerate}
    \item \label{lam:syn:fv:lem:fv-abs} $y \neq x$ असेल, तर $y \in
      \FV{\lambd[x][N]}$ जर आणि फक्त जर $y \in \FV{N}$.
    \item \label{lam:syn:fv:lem:fv-app} $y \in \FV{PQ}$ जर आणि
      फक्त जर $y \in \FV{P}$ किंवा $y
      \in \FV{Q}$.
    \end{enumerate}
\end{lem}

\begin{proof}
  सरावासाठी.
\end{proof}

\begin{prob}
  \ref{lam:syn:fv:lem:fv} सिद्ध करा.
\end{prob}











\section{आदेशन}\label{lam:syn:sub:sec}

\begin{explain}
मुक्त चरे हे परिसरातील चरांचे संदर्भ असतात. म्हणून मुक्त
चराच्या जागी ठराविक मूल्य वापरणे अर्थपूर्ण ठरते. उदाहरणार्थ,
$\lambd[x][f x]$ मधील $f$ च्या जागी ओळख फलनासारखे
$\lambd[y][y]$ हे विशिष्ट पद ठेवायचे असू शकते. त्यातून
$\lambd[x][(\lambd[y][y]) x]$ मिळते. मुक्त चरांच्या जागी
लॅम्डा पदे ठेवण्याच्या प्रक्रियेला आदेशन म्हणतात.
\end{explain}

\begin{defn}[आदेशन] \label{lam:syn:sub:defn:substitution}
  पद $M$ मधील चर $x$ च्या जागी पद $N$ चे
  \emph{आदेशन}, $\Subst{M}{N}{x}$, विगमनाने पुढीलप्रमाणे
  परिभाषित करतात:
  \begin{enumerate}
    \item $\Subst{x}{N}{x} = N$. \label{lam:syn:sub:defn:substitution-1}
    \item $\Subst{y}{N}{x}  = y$, जर $x \neq y$. \label{lam:syn:sub:defn:substitution-2}
    \item $\Subst{PQ}{N}{x}  = (\Subst{P}{N}{x}) (\Subst{Q}{N}{x})$.
      \label{lam:syn:sub:def:substitution-3}
    \item $\Subst{(\lambd[y][P])}{N}{x} = \lambd[y][\Subst{P}{N}{x}]$,
      जर $x \neq y$ आणि $y \notin \FV{N}$; अन्यथा अपरिभाषित.
      \label{lam:syn:sub:defn:substitution-4}
  \end{enumerate}
\end{defn}

\begin{explain}
\ref{lam:syn:sub:defn:substitution}\ref{lam:syn:sub:defn:substitution-4} मध्ये
$x \neq y$ ही अट घालण्याचे कारण असे की~$M$ मधील
चर~$x$ च्या \emph{बद्ध} आस्थितींच्या जागी~$N$ ठेवायचे नाही.
उदाहरणार्थ, $\Subst{\lambd[x][x]}{y}{x}$ विचारात घेताना
$\lambd[x][y]$ मिळू नये; बद्ध आस्थिती बदलली नाही तर
$\lambd[x][x]$ हेच पद राहते. मात्र वरच्या औपचारिक
नियमांनुसार या प्रसंगी आदेशन अपरिभाषित आहे.

$\lambd[y][P]$ मध्ये~$x$ च्या जागी $N$ चे आदेशन करताना
$y \notin \FV{N}$ हीही अट घालतात. उदाहरणार्थ,
$\lambd[y][x]$ मध्ये $x$ च्या जागी $y$ चे आदेशन,
म्हणजे $\Subst{\lambd[y][x]}{y}{x}$, करता येत नाही:
त्यातून $\lambd[y][y]$ मिळेल. हे पद फलसाधक स्वीकारून
तोच परत करणारे फलन दर्शवते. पण $\lambd[y][x]$ हे पद
नेहमी~$x$ (किंवा $x$ ज्याचा संदर्भ देते ते) परत करणारे
फलन दर्शवते. म्हणून अपेक्षित परिणाम म्हणजे फलसाधक
स्वीकारून तो बाजूला ठेवणारे आणि परिसरातील चर~$y$
परत करणारे फलन. हे नीट साधण्यासाठी आधी बद्ध चर~$y$
चे नाव ``बदलावे'' लागेल.
\end{explain}

\begin{prob}
  पुढील आदेशनांचे परिणाम काय आहेत?
  \begin{enumerate}
  \item $\Subst{\lambd[y][x(\lambd[w][vwx])]}{(uv)}{x}$
  \item $\Subst{\lambd[y][x(\lambd[x][x])]}{(\lambd[y][xy])}{x}$
  \item $\Subst{y(\lambd[v][xv])}{(\lambd[y][vy])}{x}$
  \end{enumerate}
\end{prob}

\begin{thm} \label{lam:syn:sub:thm:notinfv}
  $x \notin \FV{M}$ असेल, तर डावी बाजू परिभाषित
  असल्यास $\FV{\Subst{M}{N}{x}} = \FV{M}$.
\end{thm}

\begin{proof}
  $M$ च्या रचनेवर विगमन करू.
  \begin{enumerate}
  \item $M$ हे चर आहे: सरावासाठी.
  \item $M$ हे $(PQ)$ या रूपाचे आहे: सरावासाठी.
  \item $M$ हे $\lambd[y][P]$ या रूपाचे आहे. तसेच
    $\Subst{\lambd[y][P]}{N}{x}$ परिभाषित असल्यामुळे
    ते $\lambd[y][\Subst{P}{N}{x}]$ असते. म्हणून
    $\Subst{P}{N}{x}$ परिभाषित असते; शिवाय $x \neq y$
    आणि $x \notin \FV{P}$. त्यामुळे:
    \begin{multline*}
      \FV{\Subst{\lambd[y][P]}{N}{x}} = \\
    \begin{aligned}
      & = \FV{\lambd[y][\Subst{P}{N}{x}]} &&
      \text{\ref{lam:syn:sub:defn:substitution-4} नुसार}\\
      & = \FV{\Subst{P}{N}{x}} \setminus \{y\} &&
      \text{\ref{lam:syn:fv:def:fv}\ref{lam:syn:fv:def:fv2} नुसार}\\
      & = \FV{P} \setminus \{y\} && \text{विगमन गृहीतकानुसार} \\
      & = \FV{\lambd[y][P]} && \text{\ref{lam:syn:fv:def:fv}\ref{lam:syn:fv:def:fv2} नुसार}
    \end{aligned}
    \end{multline*}
  \end{enumerate}
\end{proof}

\begin{prob}
  \ref{lam:syn:sub:thm:notinfv} ची सिद्धता पूर्ण करा.
\end{prob}

\begin{thm} \label{lam:syn:sub:thm:infv}
  $x \in \FV{M}$ असेल, तर डावी बाजू परिभाषित
  असल्यास $\FV{\Subst{M}{N}{x}} = (\FV{M} \setminus
  \{x\}) \cup \FV{N}$.
\end{thm}

\begin{proof}
  $M$ च्या रचनेवर विगमन करू.
  \begin{enumerate}
  \item $M$ हे चर आहे: सरावासाठी.
  \item $M$ हे $PQ$ या रूपाचे आहे: $\Subst{(PQ)}{N}{x}$
    परिभाषित असल्यामुळे ते
    $(\Subst{P}{N}{x})(\Subst{Q}{N}{x})$ असते आणि दोन्ही
    आदेशन परिभाषित असतात. शिवाय $x \in \FV{PQ}$
    असल्यामुळे $x \in \FV{P}$ किंवा $x \in \FV{Q}$,
    किंवा दोन्ही, असते. उरलेला भाग सरावासाठी आहे.
  \item $M$ हे $\lambd[y][P]$ या रूपाचे आहे.
    $\Subst{\lambd[y][P]}{N}{x}$ परिभाषित असल्यामुळे
    ते $\lambd[y][\Subst{P}{N}{x}]$ असते;
    $\Subst{P}{N}{x}$ परिभाषित असते, $x \neq y$ आणि
    $y \notin \FV{N}$. तसेच $x \in
    \FV{\lambd[y][P]}$ असल्यामुळे $x \in \FV{P}$.
    आता:
    \begin{multline*}
      \FV{\Subst{(\lambd[y][P])}{N}{x}} = \\
      \begin{aligned}
      & = \FV{\lambd[y][\Subst{P}{N}{x}]} \\
      & = \FV{\Subst{P}{N}{x}} \setminus \{y\} \\
      & = ((\FV{P} \setminus \{x\}) \cup \FV{N}) \setminus \{y\}
       && \text{विगमन गृहीतकानुसार} \\
      & = (\FV{P} \setminus \{x, y\}) \cup \FV{N}
       && y \notin \FV{N} \\
      & = (\FV{\lambd[y][P]} \setminus \{x\}) \cup \FV{N}
      \end{aligned}
      \end{multline*}
  \end{enumerate}
\end{proof}

\begin{prob}
  \ref{lam:syn:sub:thm:infv} ची सिद्धता पूर्ण करा.
\end{prob}

\begin{thm}\label{lam:syn:sub:thm:clr}
  $x \notin \FV{N}$ आणि आदेशन परिभाषित
  असेल, तर $x \notin \FV{\Subst{M}{N}{x}}$.
\end{thm}

\begin{proof}
  सरावासाठी.
\end{proof}

\begin{prob}
  \ref{lam:syn:sub:thm:clr} सिद्ध करा.
\end{prob}


\begin{thm}\label{lam:syn:sub:thm:inv}
  $\Subst{M}{y}{x}$ परिभाषित असेल आणि $y \notin \FV{M}$,
  तर $\Subst{\Subst{M}{y}{x}}{x}{y} = M$.
\end{thm}

\begin{proof}
  $M$ च्या रचनेवर विगमन करू.
  \begin{enumerate}
  \item $M$ हे चर $z$ आहे: सरावासाठी.
  \item $M$ हे $(PQ)$ या रूपाचे आहे. मग:
    \begin{align*}
      \Subst{\Subst{(PQ)}{y}{x}}{x}{y}
      &=\Subst{((\Subst{P}{y}{x})(\Subst{Q}{y}{x}))}{x}{y} \\
      &= (\Subst{\Subst{P}{y}{x}}{x}{y})(\Subst{\Subst{Q}{y}{x}}{x}{y}) \\
      &= (PQ) \text{ विगमन गृहीतकानुसार}
    \end{align*}
  \item $M$ हे $\lambd[z][N]$ या रूपाचे आहे.
    $\Subst{\lambd[z][N]}{y}{x}$ परिभाषित असल्यामुळे
    $z \neq y$ हे माहीत आहे. म्हणून:
    \begin{align*}
      \Subst{\Subst{(\lambd[z][N])}{y}{x}}{x}{y}\\
      & = \Subst{(\lambd[z][\Subst{N}{y}{x}])}{x}{y} \\
      & = \lambd[z][\Subst{\Subst{N}{y}{x}}{x}{y}] \\
      &= \lambd[z][N] \text{ विगमन गृहीतकानुसार}
    \end{align*}
  \end{enumerate}
\end{proof}

\begin{prob}
  \ref{lam:syn:sub:thm:inv} ची सिद्धता पूर्ण करा.
\end{prob}











\section{$\alpha$-परिवर्तन}\label{lam:syn:alp:sec}

$\lambd[x][x]$ आणि $\lambd[y][y]$ यांचा संबंध काय?
ही दोन्ही पदे ओळख फलन दर्शवतात. अर्थात, विन्यासाच्या
दृष्टीने ती भिन्न पदे आहेत. त्यांच्यात फक्त बद्ध चराच्या
नावाचा फरक आहे; एकातील बद्ध चराचे ``नामांतर'' केल्यावर
दुसरे पद मिळते. याला \emph{$\alpha$-परिवर्तन} म्हणतात.

\begin{defn}[बद्ध चराचे नामांतर, $\aconvone$]
  पद $M$ मध्ये $\lambd[x][N]$ ची आस्थिती असेल,
  $x \neq y$, $y \notin \FV{N}$ असेल आणि
  $\Subst{N}{y}{x}$ परिभाषित असेल, तर ती आस्थिती
  \begin{equation*}
    \lambd[y][\Subst{N}{y}{x}]
  \end{equation*}
  ने बदलून मिळणाऱ्या $M'$ ला \emph{बद्ध चराचे नामांतर}
  म्हणतात आणि $M \redone[\alpha] M'$ असे लिहितात.
\end{defn}

\begin{defn}[संबंधाची रचनांशी सुसंगती]
  पदांवरील संबंध $R$ पुढील अटी पूर्ण करत असेल, तर तो
  \emph{रचनांशी सुसंगत} म्हणतात:
  \begin{enumerate}
  \item $R N N'$ असेल, तर $R \lambd[x][N] \lambd[x][N']$.
  \item $R P P'$ असेल, तर $R (PQ) (P'Q)$.
  \item $R Q Q'$ असेल, तर $R (PQ) (PQ')$.
  \end{enumerate}
\end{defn}

आता ही व्याख्या दुसऱ्या रीतीने मांडू:
\begin{defn}[बद्ध चराचे नामांतर, $\aconvone$]
  \emph{बद्ध चराचे नामांतर} ($\redone[\alpha]$) हा पदांवरील
  सर्वांत लहान रचनांशी सुसंगत संबंध आहे जो पुढील अट पूर्ण करतो:
  \begin{align*}
    &\lambd[x][N] \redone[\alpha] \lambd[y][\Subst{N}{y}{x}] && \text{जर
      $x \neq y$, $y \notin \FV{N}$} \\
    & &&\text{आणि $\Subst{N}{y}{x}$ परिभाषित असेल}
  \end{align*}
\end{defn}

येथे ``सर्वांत लहान'' म्हणजे सुसंगती आणि अतिरिक्त अट
यांमुळे आवश्यक ठरणाऱ्या जोड्याच या संबंधात असतात;
इतर कोणत्याही नसतात. म्हणून हा संबंध पुढीलप्रमाणेही
परिभाषित करता येतो:

\begin{defn}[बद्ध चराचे नामांतर, $\aconvone$] \label{lam:syn:alp:defn:aconvone}
  \emph{बद्ध चराचे नामांतर} ($\aconvone$) विगमनाने
  पुढीलप्रमाणे परिभाषित करतात:
  \begin{enumerate}
  \item $N \aconvone N'$ असेल, तर $\lambd[x][N] \aconvone
    \lambd[x][N']$. \label{lam:syn:alp:defn:aconvone1}
  \item $P \aconvone P'$ असेल, तर $(PQ) \aconvone (P'Q)$. \label{lam:syn:alp:defn:aconvone2}
  \item $Q \aconvone Q'$ असेल, तर $(PQ) \aconvone (PQ')$. \label{lam:syn:alp:defn:aconvone3}
  \item $x \neq y$, $y \notin \FV{N}$ आणि $\Subst{N}{y}{x}$
    परिभाषित असेल, तर
    $\lambd[x][N] \redone[\alpha] \lambd[y][\Subst{N}{y}{x}]$.
    \label{lam:syn:alp:defn:aconvone4}
  \end{enumerate}
\end{defn}

या व्याख्या सममूल्य आहेत; त्यांची सिद्धता सरावासाठी ठेवतो.
यापुढे आपण विगमनाधारित व्याख्या वापरू.

\begin{defn}[$\alpha$-परिवर्तन, $\aconv$]
  \emph{$\alpha$-परिवर्तन} ($\aconv$) हा पदांवरील
  $\aconvone$ समाविष्ट करणारा सर्वांत लहान परावर्ती
  आणि संक्रमक संबंध आहे.
\end{defn}

वरच्याप्रमाणे, ``सर्वांत लहान'' म्हणजे परावर्तिता, संक्रमणशीलता आणि
$\aconvone$ यांमुळे आवश्यक ठरणाऱ्या जोड्याच या संबंधात
असतात. यावरून पुढील सममूल्य व्याख्या मिळते:

\begin{defn}[$\alpha$-परिवर्तन, $\aconv$] \label{lam:syn:alp:defn:aconv}
  \emph{$\alpha$-परिवर्तन} ($\aconv$) विगमनाने
  पुढीलप्रमाणे परिभाषित करतात:
  \begin{enumerate}
  \item $P \aconv Q$ आणि $Q \aconv R$ असतील, तर $P \aconv R$.
    \label{lam:syn:alp:defn:aconv1}
  \item $P \aconvone Q$ असेल, तर $P \aconv Q$. \label{lam:syn:alp:defn:aconv2}
  \item $P \aconv P$. \label{lam:syn:alp:defn:aconv3}
  \end{enumerate}
\end{defn}

\begin{ex}
  $\lambd[x][f x]$ चे $\alpha$-परिवर्तन करून
  $\lambd[y][f
  y]$ मिळते; उलट दिशेनेही परिवर्तन करता येते.
  अनौपचारिकपणे, ही दोन्ही पदे फलसाधक स्वीकारून
  त्यावर $f$ लावल्याचा परिणाम देणारी फलने आहेत;
  त्यांतील~$f$ हा परिसरातील चराचा संदर्भ आहे.

  $\lambd[x][f x]$ चे $\alpha$-परिवर्तन करून
  $\lambd[x][g
    x]$ मिळत नाही. ही पदे अनुक्रमे परिसरातील
  $f$ आणि~$g$ यांचा संदर्भ घेतात आणि म्हणून भिन्न
  फलने दर्शवतात: $f$ आणि $g$ भिन्न असलेल्या परिसरांत
  त्यांचे वर्तन भिन्न असते.
\end{ex}

\begin{prob}
  पुढील पदजोड्या $\alpha$-परिवर्तनीय आहेत का?
  \begin{enumerate}
  \item $\lambd[x][\lambd[y][x]]$ आणि $\lambd[y][\lambd[x][y]]$
  \item $\lambd[x][\lambd[y][x]]$ आणि $\lambd[c][\lambd[b][a]]$
  \item $\lambd[x][\lambd[y][x]]$ आणि $\lambd[c][\lambd[b][a]]$
  \end{enumerate}
\end{prob}

\begin{lem}\label{lam:syn:alp:lem:fv-one}
  $P \aconvone Q$ असेल, तर $\FV{P} = \FV{Q}$.
\end{lem}

\begin{proof}
  $P \aconvone Q$ च्या निष्पत्ती वर विगमन करू.
  \begin{enumerate}
  \item शेवटचा नियम \ref{lam:syn:alp:defn:aconvone4} असेल, तर $P$
    हे $\lambd[x][N]$ या रूपाचे आणि $Q$ हे
    $\lambd[y][\Subst{N}{y}{x}]$ या रूपाचे असते;
    येथे $x \neq y$, $y \notin \FV{N}$ आणि
    $\Subst{N}{y}{x}$ परिभाषित असते. $x \in \FV{N}$
    आहे की नाही यानुसार प्रकरणे वेगळी करू:
    \begin{enumerate}
    \item $x \in \FV{N}$ असेल, तर:
      \begin{align*}
        \FV{\lambd[y][\Subst{N}{y}{x}]} &=
        \FV{\Subst{N}{y}{x}} \setminus \{y\} \\
        & = ((\FV{N} \setminus \{x\}) \cup \{y\}) \setminus \{y\}
         && \text{\ref{lam:syn:sub:thm:infv} नुसार} \\
        & = \FV{N} \setminus \{x\} \\
        & = \FV{\lambd[x][N]}
      \end{align*}
    \item $x \notin \FV{N}$ असेल, तर:
      \begin{align*}
        \FV{\lambd[y][\Subst{N}{y}{x}]}
        & = \FV{\Subst{N}{y}{x}} \setminus \{y\} \\
        & = \FV{N} \setminus \{x\}
         && \text{\ref{lam:syn:sub:thm:notinfv} नुसार} \\
        & = \FV{\lambd[x][N]}.
      \end{align*}
    \end{enumerate}
  \item उरलेली तीन प्रकरणे सरावासाठी आहेत.
  \end{enumerate}
\end{proof}

\begin{prob}
  \ref{lam:syn:alp:lem:fv-one} ची सिद्धता पूर्ण करा.
\end{prob}

\begin{lem}\label{lam:syn:alp:lem:inv}
  $P \aconvone Q$ असेल, तर $Q \aconvone P$.
\end{lem}

\begin{proof}
  $P \aconvone Q$ च्या निष्पत्ती वर विगमन करू.
  \begin{enumerate}
  \item शेवटचा नियम \ref{lam:syn:alp:defn:aconvone4} असेल, तर $P$
    हे $\lambd[x][N]$ या रूपाचे आणि $Q$ हे
    $\lambd[y][\Subst{N}{y}{x}]$ या रूपाचे असते;
    येथे $x \neq y$, $y \notin \FV{N}$ आणि
    $\Subst{N}{y}{x}$ परिभाषित असते. प्रथम,
    \ref{lam:syn:sub:thm:clr} नुसार $x \notin
    \FV{\Subst{N}{y}{x}}$ आहे. \ref{lam:syn:sub:thm:inv}
    नुसार $\Subst{\Subst{N}{y}{x}}{x}{y}$ केवळ
    परिभाषितच नाही, तर~$N$ च्या बरोबरही आहे. म्हणून
    \ref{lam:syn:alp:defn:aconvone4} नुसार
    $\lambd[y][\Subst{N}{y}{x}]
    \aconvone \lambd[x][\Subst{\Subst{N}{y}{x}}{x}{y}] =
    \lambd[x][N]$.
  \end{enumerate}
\end{proof}

\begin{prob}
  \ref{lam:syn:alp:lem:inv} ची सिद्धता पूर्ण करा.
\end{prob}

\begin{thm}
  $\alpha$-परिवर्तन हा पदांवरील तुल्यता संबंध आहे;
  म्हणजे तो परावर्ती, सममित आणि संक्रमक आहे.
\end{thm}

\begin{proof}
  \begin{enumerate}
  \item प्रत्येक पद $M$ साठी, बद्ध चरांच्या \emph{शून्य}
    नामांतरांनी $M$ चे $M$ मध्ये परिवर्तन होते.
  \item $P$ चे बद्ध चरांच्या काही नामांतरांनी $Q$ मध्ये
    $\alpha$-परिवर्तन होत असेल, तर \ref{lam:syn:alp:lem:inv}
    वापरून ती नामांतरे उलट क्रमाने उलटवता येतात आणि
    $Q$ पासून~$P$ मिळते.
  \item $P$ चे बद्ध चरांच्या एका क्रमाने $Q$ मध्ये आणि
    $Q$ चे दुसऱ्या क्रमाने $R$ मध्ये $\alpha$-परिवर्तन
    होत असेल, तर दोन्ही क्रम सलग लावून $P$ पासून
    $R$ मिळते.
  \end{enumerate}
\end{proof}

यापुढे $M$ आणि $N$ ही पदे \emph{$\alpha$-सममूल्य},
$M \aeq N$, आहेत असे म्हणू, जर आणि फक्त जर $M$ चे $N$ मध्ये
$\alpha$-परिवर्तन होत असेल. आपण नुकतेच दाखवल्याप्रमाणे,
हे घडते, जर आणि फक्त जर $N$ चेही~$M$ मध्ये
$\alpha$-परिवर्तन होते.

\begin{thm}\label{lam:syn:alp:thm:fv}
  $M \aeq N$ असेल, तर $\FV{M} = \FV{N}$.
\end{thm}

\begin{proof}
  \ref{lam:syn:alp:lem:fv-one} वरून लगेच मिळते.
\end{proof}

\begin{lem}\label{lam:syn:alp:lem:sub:R}
  $R \aeq R'$ असेल आणि $\Subst{M}{R}{y}$ परिभाषित
  असेल, तर $\Subst{M}{R'}{y}$ परिभाषित असते आणि
  $\Subst{M}{R}{y}$ शी $\alpha$-सममूल्य असते.
\end{lem}

\begin{proof}
  सरावासाठी.
\end{proof}

\begin{prob}
  \ref{lam:syn:alp:lem:sub:R} सिद्ध करा.
\end{prob}

\ref{lam:syn:sub:sec} मध्ये पाहिल्याप्रमाणे, काही प्रसंगी
आदेशन अपरिभाषित असते. पण पदांवर $\alpha$-परिवर्तन
वापरून बद्ध चरांची नावे बदलल्यास आदेशन नेहमी परिभाषित
करता येते. परिणामी पदाची $\alpha$-सममूल्यता टिकते,
हे पुढील प्रमेय दाखवते:

\begin{thm}\label{lam:syn:alp:thm:sub}
  कोणत्याही $M$, $R$ आणि~$y$ साठी असे $M'$ असते की
  $M \aeq M'$ आणि $\Subst{M'}{R}{y}$ परिभाषित असते.
  आणखी एखादी जोडी $M'' \aeq M$ आणि $R''$ अशी असेल
  की $\Subst{M''}{R''}{y}$ परिभाषित आहे आणि $R'' \aeq R$,
  तर $\Subst{M'}{R}{y} \aeq \Subst{M''}{R''}{y}$.
\end{thm}

\begin{proof}

प्रथम एक सांत संच वर्ज्य नावे म्हणून दिला असता, कोणत्याही
पदातील सर्व बद्ध नावे त्या संचाबाहेर, परस्पर भिन्न आणि
पदाच्या सर्व मूळ नावांपासून वेगळी करता येतात, हे पाहू.
सर्वांत आतल्या अमूर्तीकरणांपासून बाहेरच्या दिशेने नामांतर
करा. प्रत्येक पायरीला आधी निवडलेल्या नावांपासूनही वेगळे
नवे नाव निवडा; अमर्याद चर असल्याने अशी निवड शक्य आहे.
या पायरीत आतल्या सर्व binder ची नावे आधीच ताजी आहेत,
म्हणून जुन्या बाह्य binder चे नाव त्यांच्याशी समान नाही.
नवे नावही त्यांच्याशी किंवा देहातील मुक्त नावांशी समान नाही.
त्यामुळे \ref{lam:syn:sub:defn:substitution} मधील आंशिक
नियमांनी ही renaming substitution परिभाषित असते आणि
\ref{lam:syn:alp:defn:aconvone4} लागू करता येतो. पदाच्या रचनेवरील
विगमनाने प्रत्येक पायरी आणि संपूर्ण सांत प्रक्रिया योग्य ठरते.
आता वर्ज्य नावांत $y$ आणि $\FV{R}$ मधील सर्व नावे घ्या
आणि ही प्रक्रिया $M$ वर करा; मिळालेल्या पदाला $M'$ म्हणा.
प्रत्येक पायरी alpha परिवर्तन असल्याने $M\aeq M'$.
आता कोणतेही बद्ध नाव $y$ नाही किंवा $R$ मध्ये मुक्त नाही.
म्हणून रचनेवरील सरळ विगमनाने $\Subst{M'}{R}{y}$
परिभाषित आहे. याने अस्तित्व सिद्ध होते.

एकमेव alpha-class मिळण्याचा उरलेला दावा पुढीलप्रमाणे
सिद्ध करू. \ref{lam:syn:unq:prop:unq} मधील एकमेव रचनावृक्षात
प्रत्येक मुक्त चराच्या पानावर त्याचे नाव ठेवा; प्रत्येक बद्ध
वापराला त्याला बांधणाऱ्या अमूर्तीकरणाच्या नोडचा निर्देश द्या.
अमूर्तीकरणांच्या नावांखेरीज हा वृक्ष आणि हे निर्देश प्रत्येक
alpha पायरीत जपले जातात. हे rename नियमाच्या ताजेपणाच्या
अटीवरून आणि तीन compatibility नियमांवरून तपासता येते.
उलट, दोन पदांचे असे नामनिरपेक्ष वृक्ष व निर्देश समान असतील,
तर संबंधित अमूर्तीकरणांच्या प्रत्येक स्थानाला दोन्ही पदांच्या
सर्व नावांबाहेरचे तेच ताजे नाव द्या. वरील आतून-बाहेर
नामांतरप्रक्रियेने दोन्ही पदे एकाच नामांकित पदात बदलतात.
म्हणून ती alpha-सममूल्य आहेत.
जेव्हा raw आदेशन परिभाषित असते, तेव्हा ते फक्त मुक्त
चराच्या संबंधित पानांच्या जागी आदेशित पदाचा असा वृक्ष बसवते.
आंशिक नियमातील अटींमुळे त्यातील मुक्त पानांना बाह्य binder
पकडत नाहीत; मूळ बद्ध पानांचे निर्देशही बदलत नाहीत.
त्यामुळे $M'\aeq M''$ आणि $R\aeq R''$ असताना,
दोन्ही आदेशन परिभाषित असतील तर त्यांचे नामनिरपेक्ष
परिणामी वृक्ष समान असतात. वरील उलट दिशेने
$\Subst{M'}{R}{y}\aeq\Subst{M''}{R''}{y}$ मिळते.
याने अस्तित्व आणि परिणामी alpha-class ची uniqueness
दोन्ही दावे सिद्ध होतात.
\end{proof}

\begin{prob}
  \ref{lam:syn:alp:thm:sub} ची सिद्धता पूर्ण करा.
\end{prob}

\begin{cor}\label{lam:syn:alp:cor:sub}
  कोणत्याही $M$, $R$ आणि~$y$ साठी $M'$ आणि $R'$ ची
  अशी जोडी असते की $M' \aeq M$, $R \aeq R'$
  आणि $\Subst{M'}{R'}{y}$ परिभाषित असते. आणखी
  $M'' \aeq M$ आणि $R'' \aeq R$ अशी जोडी असेल
  आणि $\Subst{M''}{R''}{y}$ परिभाषित असेल, तर
  $\Subst{M'}{R'}{y} \aeq
  \Subst{M''}{R''}{y}$.
\end{cor}

\begin{proof}
  \ref{lam:syn:alp:thm:sub} वरून लगेच मिळते.
\end{proof}












\section{डी ब्रुइन निर्देशांक}\label{lam:syn:deb:sec}

$\alpha$-सममूल्यता ही स्वाभाविक कल्पना आहे, कारण
$\alpha$-सममूल्य पदांचा ``अर्थ एकच'' असतो. प्रत्यक्षात,
एकमेकांत $\alpha$-परिवर्तन करता येणारी पदे वेगळी
न मानणारा लॅम्डा पदांचा विन्यास देता येतो. यांपैकी
सर्वाधिक परिचित पद्धतीत चरांच्या जागी त्यांचे
\emph{डी ब्रुइन निर्देशांक} वापरतात.

$\lambd[x][M]$ लिहिताना फलनाचा प्राचल $x$ आहे हे
स्पष्ट सांगतो, म्हणजे~$M$ मध्ये त्या प्राचलाचा संदर्भ
घेण्यासाठी $x$ वापरता येतो. मात्र डी ब्रुइन निर्देशांक
पद्धतीत प्राचलांना नावे नसतात. फलनाच्या मुख्य भागात
प्राचलाचा संदर्भ हा त्याच्या आस्थितीपासून त्याला बद्ध
करणाऱ्या अमूर्तीकरणापर्यंत मधे येणाऱ्या अमूर्तीकरणांची
पातळी मोजणाऱ्या संख्येने दर्शवतात. उदाहरणार्थ,
$\lambd[x][\lambd[y][y x]]$ पाहा: बाहेरील अमूर्तीकरण
चर~$x$ बद्ध करते; आतील अमूर्तीकरण चर~$y$ बद्ध करते.
उपपद $y x$ हे आतील अमूर्तीकरणाच्या व्याप्तीत आहे.
$y$ आणि त्याला बद्ध करणारे~$\lambd[y]$ यांच्यात
एकही अमूर्तीकरण नाही; पण $x$ आणि त्याला बद्ध
करणारे~$\lambd[x]$ यांच्यात एक अमूर्तीकरण आहे.
म्हणून $y x$ साठी $0\, 1$, आणि संपूर्ण पदासाठी
$\lambd[][\lambd[][01]]$ असे लिहितात.

\begin{defn}
  डी ब्रुइन पदे विगमनाने पुढीलप्रमाणे परिभाषित करतात:
  \begin{enumerate}
  \item $n$, जेथे $n$ ही कोणतीही नैसर्गिक संख्या आहे.
  \item $PQ$, जेथे $P$ आणि $Q$ ही दोन्ही डी ब्रुइन पदे आहेत.
  \item $\lambd[][N]$, जेथे $N$ हे डी ब्रुइन पद आहे.
  \end{enumerate}
\end{defn}


सामान्य लॅम्डा पदांचे डी ब्रुइन निर्देशांकित पदांत
औपचारिक रूपांतर पुढीलप्रमाणे करतात:
\begin{defn}
  \begin{align*}
    F_\Gamma(x) &= \Gamma(x) \\
    F_\Gamma(PQ) &= F_\Gamma(P)F_\Gamma(Q) \\
    F_\Gamma(\lambd[x][N]) &= \lambd[][F_{x,\Gamma}(N)]
  \end{align*}
  येथे $\Gamma$ ही शून्यापासून निर्देशांकित चरांची यादी
  आहे आणि $\Gamma(x)$ ही $\Gamma$ मधील चर $x$ ची
  डावीकडून पहिली जागा दर्शवते. त्यामुळे एकाच नावाची
  पुनरावृत्ती असली, तरी सर्वांत जवळच्या अमूर्तीकरणाचा
  निर्देशांक मिळतो. सुरुवातीच्या यादीत रूपांतर करायच्या
  पदाची सर्व मुक्त चरे असणे आवश्यक आहे. उदाहरणार्थ, $\Gamma$ ही $x,y,z$
  असेल, तर $\Gamma(x)$ हे $0$ आणि $\Gamma(z)$ हे $2$ आहे.

  $x,\Gamma$ म्हणजे $\Gamma$ च्या सुरुवातीला $x$
  ठेवून मिळालेली यादी. उदाहरणार्थ, मागील उदाहरणातील
  $\Gamma$ साठी $w,\Gamma$ ही $w,x,y,z$ आहे.
\end{defn}

डी ब्रुइन पदापासून नेहमीचे लॅम्डा पद पुढीलप्रमाणे
परत मिळवतात. हे रूपांतर तेव्हाच परिभाषित असते, जेव्हा
प्रत्येक संख्या त्या आस्थितीवर उपलब्ध असलेल्या यादीतील
एका जागेचा निर्देशांक असते. अमूर्तीकरणाच्या देहात
जाताना यादीच्या सुरुवातीला आणखी एक नाव येते.
उपलब्ध जागेबाहेरचा निर्देशांक असलेल्या पदासाठी
हे उलट रूपांतर परिभाषित नाही:

\begin{defn}
  \begin{align*}
    G_\Gamma(n) &= \Gamma[n] \\
    G_\Gamma(PQ) &= G_\Gamma(P) G_\Gamma(Q) \\
    G_\Gamma(\lambd[][N]) &= \lambd[x][G_{x,\Gamma}(N)]
  \end{align*}
  येथेही $\Gamma$ ही शून्यापासून निर्देशांकित चरांची
  यादी आहे, आणि $\Gamma[n]$ म्हणजे~$n$ व्या जागेवरील
  चर. उदाहरणार्थ, $\Gamma$ ही $x,y,z$ असेल, तर
  $\Gamma[1]$ हे $y$ आहे.

  शेवटच्या समीकरणातील चर $x$ हा $\Gamma$ मध्ये
  नसलेला कोणताही चर निवडतात.
\end{defn}

येथे काही निष्कर्ष त्यांच्या सिद्धता न देता मांडतो:

\begin{prop}
  $M \aconvone M'$ असेल, आणि $\Gamma$ ही
  $\FV{M}$ समाविष्ट करणारी कोणतीही यादी असेल,
  तर $F_\Gamma(M) \eqs F_\Gamma(M')$.
\end{prop}











\section{$\alpha$-सममूल्यतावर्ग म्हणून पदे}\label{lam:syn:tr:sec}

यापुढे आपण पदांचा विचार $\alpha$-सममूल्यतेपर्यंत करू.
म्हणजे एखादे पद लिहिले की ते ज्या $\alpha$-सममूल्यतावर्गात
आहे तो वर्ग अभिप्रेत असेल. उदाहरणार्थ,
$\lambd[a][\lambd[b][a c]]$ असे लिहिल्यावर त्याच्याशी
$\alpha$-सममूल्य असलेल्या सर्व पदांचा संच अभिप्रेत आहे;
त्यात $\lambd[a][\lambd[b][a
    c]]$, $\lambd[b][\lambd[a][b c]]$ इत्यादी येतात.

तसेच मागील विभागांत $N, Q$ अशी अक्षरे एखादे पद
दर्शवत होती; यापुढे ती एक वर्ग दर्शवतील. पुढील अभ्यासाचे
विषय स्वतंत्र पदांऐवजी हे वर्ग असतील. $x, y$ सारखी
अक्षरे मात्र चरच दर्शवतील.

वर्ग~$M$ मधील कोणताही घटक दर्शवण्यासाठी
$\rep{M}$ हे संकेतन वापरू. एकाहून अधिक घटक हवे
असतील, तर $\rep{M}[0], \rep{M}[1]$ इत्यादी असे लिहू.

मांडणी सोपी करण्यासाठी पदांसाठीची संकेते वर्गांसाठीही
वापरू. वर्गांवर पुढील व्याख्या करतो:
\begin{defn}
  \begin{enumerate}
  \item $\lambd[x][N]$ म्हणजे $\lambd[x][\rep{N}]$
    समाविष्ट करणारा वर्ग.
  \item $PQ$ म्हणजे $\rep{P}\rep{Q}$ समाविष्ट करणारा वर्ग.
  \end{enumerate}
\end{defn}

या व्याख्या नीट परिभाषित आहेत हे सहज दिसते,
कारण $\alpha$-परिवर्तन रचनांशी सुसंगत आहे.

\begin{defn} \label{lam:syn:tr:def:fv}
  $\alpha$-सममूल्यतावर्ग~$M$ ची \emph{मुक्त चरे},
  म्हणजे $FV(M)$, ही $FV(\rep{M})$ अशी परिभाषित करतो.
\end{defn}

ही व्याख्या नीट परिभाषित आहे, कारण
$FV(\rep{M}[0]) =
FV(\rep{M}[1])$ हे \ref{lam:syn:alp:thm:fv} मध्ये दाखवले आहे.

वर्गांमध्ये आदेशन करण्यासाठीही पूर्वीचे संकेतन वापरू:
\begin{defn} \label{lam:syn:tr:def:sub}
  $M$ मध्ये $y$ च्या जागी $R$ चे \emph{आदेशन}, म्हणजे
  $\Subst{M}{R}{y}$, हा $\Subst{\rep{M}}{\rep{R}}{y}$
  समाविष्ट करणारा वर्ग म्हणून परिभाषित करतो; येथे
  $\rep{M}$ आणि $\rep{R}$ असे निवडायचे की आदेशन
  परिभाषित असेल.
\end{defn}

\ref{lam:syn:alp:cor:sub} नुसार ही व्याख्याही नीट परिभाषित आहे.

या व्याख्येमुळे युक्तिवाद किती सोपा होतो ते पाहा.
उदाहरणार्थ, पुढील z हे आदेशित चर आणि त्याच्या जागी
बसवायचे चर या दोन्हींपासून वेगळे ताजे नाव निवडा:
\begin{align}
  \Subst{\lambd[x][x]}{y}{x} & =\label{lam:syn:tr:eq:1}\\
  &= \Subst{\lambd[z][z]}{y}{x} \label{lam:syn:tr:eq:2}\\
  &= \lambd[z][\Subst{z}{y}{x}] \\
  &= \lambd[z][z]
\end{align}

\ref{lam:syn:tr:eq:1} कडे अजूनही स्वतंत्र पदांवरील आदेशन म्हणून
पाहिले, तर ते अपरिभाषित आहे. पण आधी सांगितल्याप्रमाणे,
आता आपण वर्गांवरील आदेशन विचारात घेतो; म्हणून
\ref{lam:syn:tr:eq:2} करता येते. $\lambd[x][x]$ आणि
$\lambd[z][z]$ हे एकाच वर्गात असल्याने पहिल्या
पदाऐवजी दुसरे वापरता येते.

याच कारणाने यापुढे आदेशनासाठी आवश्यक अटी पूर्ण
करणारेच प्रतिनिधी आपण निवडले आहेत असे गृहीत धरू.
उदाहरणार्थ, $\Subst{\lambd[x][N]}{R}{y}$ दिसल्यास
त्यातील प्रतिनिधी $\lambd[x][N]$ असा निवडला आहे की
$x \neq y$ आणि $x \notin FV(R)$, असे मानू.
देहातील सर्व आतल्या बद्ध नावांनाही आदेशनाच्या चरापासून
आणि आदेशित पदाच्या मुक्त चरांपासून वेगळे निवडायचे आहे;
केवळ बाहेरच्या अमूर्तीकरणाच्या नावाच्या अटी पुरेशा नाहीत.

$\lambd[x][x]$ ला ``वर्ग'' म्हणणे थोडे विचित्र वाटते.
म्हणून यापुढे अशा वर्गांना $\Lambda$-पदे (लॅम्डा-वर्गपदे)
म्हणू; या भागात पुढे त्यांना साधे ``पदे'' असेही म्हणू.
यामुळे आधीची परिचित $\lambda$-पदे आणि हे वर्ग
यांतील फरक स्पष्ट राहील.

\begin{editorial}
  मात्र पूर्वीच्या पदांचा विचार अजून सोडता येत नाही:
  $\Lambda$-पदांची संपूर्ण व्याख्या $\lambda$-पदांवर
  आधारलेली आहे. तसेच $\Lambda$-पदांवरील फलने थेट
  परिभाषित करण्याची पद्धत आपण अजून दिलेली नाही.
  त्यामुळे अशी फलने प्रथम $\lambda$-पदांवर परिभाषित
  करून मग त्यांना $\Lambda$-पदांवर ``प्रक्षेपित'' करावे
  लागते; आदेशनासाठी आपण तेच केले. मात्र वाचकाला
  $\Lambda$-पदांवरील फलने कशी परिभाषित करता येतील
  हे सहज समजेल असे आपण गृहीत धरतो.
\end{editorial}










\section{$\beta$-न्यूनीकरण}\label{lam:int:bet:sec}

$(\lambd[m][(\lambd[y][y]) m])$ पाहिल्यावर त्याचा
$\lambd[m][m]$ शी काही संबंध असावा असे वाटते:
दुसरे पद पहिल्याचे ``सोपीकरण'' केल्यावर मिळायला हवे.
\emph{$\beta$-न्यूनीकरण} ही कल्पना औपचारिकपणे पकडते.

\begin{defn}[$\beta$-संकुचन, $\bredone$] \label{lam:int:bet:defn:betacontr}
  \emph{$\beta$-संकुचन} ($\bredone$) हा पदांवरील
  सर्वांत लहान रचनांशी सुसंगत संबंध आहे जो पुढील अट पूर्ण करतो:
  \[    (\lambd[x][N])Q \bredone \Subst{N}{Q}{x}
  \]
  $P \bredone Q$ असेल, तर $P$ चे $Q$ मध्ये
  \emph{$\beta$-संकुचन} झाले असे म्हणतो.
  $(\lambd[x][N])Q$ या रूपाच्या पदाला \emph{रेडेक्स} म्हणतात.
\end{defn}

\begin{prob} \label{lam:int:bet:prob:def}
  \ref{lam:syn:alp:defn:aconvone} मध्ये बद्ध चराच्या
  नामांतरासाठी जशा मांडल्या, तशाच $\beta$-संकुचनाच्या
  सममूल्य विगमनाधारित व्याख्या सविस्तर लिहा.
\end{prob}

\begin{defn}[$\beta$-न्यूनीकरण, $\bred$] \label{lam:int:bet:defn:betared}
  \emph{$\beta$-न्यूनीकरण} ($\bred$) हा $\bredone$
  समाविष्ट करणारा पदांवरील सर्वांत लहान परावर्ती आणि
  संक्रमक संबंध आहे. $P \bred Q$ असेल, तर $P$ चे
  $Q$ मध्ये $\beta$-न्यूनीकरण झाले असे म्हणतो.
\end{defn}

संदर्भ स्पष्ट असेल, तर $\bredone$ ऐवजी $\redone$
आणि $\bred$ ऐवजी $\red$ असे लिहू.

अनौपचारिकपणे, $M \bred N$ असते, जर आणि फक्त जर
$\beta$-संकुचनाच्या शून्य किंवा अधिक पायऱ्यांनी
$M$ चे $N$ मध्ये परिवर्तन करता येते.

\begin{defn}[$\beta$-सामान्य]
ज्या पदाचे आणखी $\beta$-संकुचन करता येत नाही,
त्याला \emph{$\beta$-सामान्य} म्हणतात.
\end{defn}

$M \bred N$ आणि $N$ हे $\beta$-सामान्य असेल, तर
$N$ ला~$M$ चे \emph{सामान्य रूप} म्हणतो. एखाद्या पदाचे
सामान्य रूप एकमेव असते का, असा प्रश्न पडू शकतो.
उत्तर होय आहे, हे आपण पुढे पाहू.

काही उदाहरणे पाहू.
\begin{enumerate}
\item आपल्याला मिळते:
\begin{align*}
(\lambd[x][xxy]) \lambd[z][z] & \redone (\lambd[z][z])(\lambd[z][z]) y \\
& \redone (\lambd[z][z]) y \\
& \redone y
\end{align*}
\item पदाचे ``सोपीकरण'' केल्यावर ते प्रत्यक्षात अधिक गुंतागुंतीचेही होऊ शकते:
\begin{align*}
(\lambd[x][xxy])(\lambd[x][xxy]) & \redone (\lambd[x][xxy])(\lambd[x][xxy])y \\
& \redone (\lambd[x][xxy])(\lambd[x][xxy])yy \\
& \redone \dots
\end{align*}
\item पद तसेचही राहू शकते:
\[(\lambd[x][xx])(\lambd[x][xx]) \redone (\lambd[x][xx])(\lambd[x][xx])
\]
\item तसेच काही पदांचे एकाहून अधिक मार्गांनी न्यूनीकरण
  करता येते. उदाहरणार्थ,
\[(\lambd[x][(\lambd[y][yx]) z]) v \redone (\lambd[y][yv]) z
\]
हे बाहेरील उपयोजनाचे संकुचन करून मिळते; आणि
\[(\lambd[x][(\lambd[y][yx]) z]) v \redone (\lambd[x][zx]) v
\]
हे आतील उपयोजनाचे संकुचन करून मिळते. तथापि,
या प्रसंगी दोन्ही पदांचे पुढील न्यूनीकरण एकाच
$zv$ या पदात होते.
\end{enumerate}

शेवटच्या उदाहरणातील समान परिणाम हा योगायोग नाही.
तो लॅम्डा कलनाचा चर्च–रॉसर गुणधर्म म्हणून ओळखला
जाणारा खोल आणि महत्त्वाचा गुणधर्म दाखवतो.

\begin{digress}
  सर्वसाधारणपणे, एखाद्या पदाचे $\beta$-न्यूनीकरण
  अनेक प्रकारे करता येते. म्हणून न्यूनीकरणाच्या अनेक
  पद्धती मांडल्या गेल्या आहेत. त्यांपैकी सर्वाधिक प्रचलित
  \emph{स्वाभाविक पद्धत} आहे. ती नेहमी \emph{डावीकडील
  सर्वप्रथम} रेडेक्सचे संकुचन करते; रेडेक्सची जागा
  पदात तो जिथे सुरू होतो त्यावरून ठरवली जाते.
  या पद्धतीचा उपयुक्त गुणधर्म असा: एखाद्या पदाचे
  कोणत्याही पद्धतीने सामान्य रूपात न्यूनीकरण करता येते
  तेव्हाच स्वाभाविक पद्धतीनेही ते करता येते. पुढे
  वेगळे स्पष्ट केलेले नसेल, तर आपण स्वाभाविक पद्धत वापरू.
\end{digress}

\begin{defn}[$\beta$-सममूल्यता, $\equal$]
  \emph{$\beta$-सममूल्यता} ($\equal$) हा संबंध विगमनाने
  पुढीलप्रमाणे परिभाषित करतात:
  \begin{enumerate}
  \item $M \equal M$.
  \item $M \equal N$ असेल, तर $N \equal M$.
  \item $M \equal N$, $N \equal O$ असतील, तर $M \equal O$.
  \item $M \equal N$ असेल, तर $PM \equal PN$.
  \item $M \equal N$ असेल, तर $MQ \equal NQ$.
  \item $M \equal N$ असेल, तर $\lambd[x][M] \equal \lambd[x][N]$.
  \item $(\lambd[x][N])Q \equal \Subst{N}{Q}{x}$.
  \end{enumerate}
\end{defn}

पहिले तीन नियम या संबंधाला तुल्यता संबंध बनवतात;
पुढील तीन त्याला रचनांशी सुसंगत बनवतात; आणि शेवटचा
नियम त्यात $\beta$-संकुचन समाविष्ट करतो.

अनौपचारिकपणे, दोन पदे $\beta$-सममूल्य असतात, जर आणि फक्त जर
एका पदाचे दुसऱ्यात शून्य किंवा अधिक पायऱ्यांनी
परिवर्तन करता येते, आणि प्रत्येक पायरी $\beta$-संकुचनाची
किंवा $\beta$-संकुचनाच्या ``उलट'' प्रकारची असते.
उलट $\beta$-संकुचन
असे परिभाषित करतात की $M$ चे उलट $\beta$-संकुचन
 $N$ मध्ये होते, जर आणि फक्त जर $N$ चे $\beta$-संकुचन $M$
मध्ये होते.

वरील नियमांखेरीज आणखी नियम जोडून हा संबंध विस्तारू;
विस्तारित सममूल्यता संबंध $\equal[X]$ असा दर्शवू,
जेथे $X$ हा जोडलेला नियम आहे.











\section{$\eta$-परिवर्तन}\label{lam:syn:eta:sec}

$\lambda$-पदांवर आणखी एक संबंध आहे. \ref{lam:syn:fv:sec}
मध्ये आपण $\lambd[x][(fx)]$ हे उदाहरण पाहिले: ते
फलसाधक स्वीकारून त्यावर $f$ लावते. दुसऱ्या शब्दांत,
ते~$f$ हेच फलन दर्शवते: $\lambd[x][(fx)]N$ आणि
$fN$ या दोन्हींचे न्यूनीकरण~$fN$ मध्ये होते. ही कल्पना
मांडण्यासाठी $\eta$-न्यूनीकरण (आणि $\eta$-विस्तार)
वापरतो.

\begin{defn}[$\eta$-संकुचन, $\eredone$]
  \label{lam:syn:eta:defn:beredone}
  \emph{$\eta$-संकुचन} ($\eredone$) हा पदांवरील
  सर्वांत लहान रचनांशी सुसंगत संबंध आहे जो पुढील अट पूर्ण करतो:
  \[    \lambd[x][M x] \eredone M \text{ जर } x \notin FV(M)
  \]
\end{defn}

\begin{defn}[$\beta\eta$-न्यूनीकरण, $\bered$] \label{lam:syn:eta:defn:bered}
  \emph{$\beta\eta$-न्यूनीकरण} ($\bered$) हा $\bredone$
  आणि $\eredone$ समाविष्ट करणारा पदांवरील सर्वांत लहान
  परावर्ती आणि संक्रमक संबंध आहे. म्हणजे परावर्तन आणि
  संक्रमणाचे नियम, तसेच पुढील दोन नियम, त्यात असतात:
  \begin{enumerate}
  \item $M \bredone N$ असेल, तर $M \bered N$. \label{lam:syn:eta:defn:bered3}
  \item $M \eredone N$ असेल, तर $M \bered N$. \label{lam:syn:eta:defn:bered4}
  \end{enumerate}

\end{defn}

\begin{defn}
  $\equal$ हा सममूल्यता संबंध $\eta$-परिवर्तनाच्या
  पुढील नियमाने विस्तारतो:
  \[  \lambd[x][M x] \equal M \text{ जर } x \notin FV(M)
  \]
  विस्तारित संबंध $\equal[\eta]$ असा दर्शवतो.
\end{defn}

$\eta$-सममूल्यता महत्त्वाची आहे, कारण तिचा लॅम्डा
पदांच्या विस्तारात्मकतेशी संबंध आहे:

\begin{defn}[विस्तारात्मकता]
  $\equal$ हा सममूल्यता संबंध (\ext) नियमाने विस्तारतो:
  \begin{center}
  $Mx \equal Nx$ असेल, तर $M \equal N$; येथे $x \notin FV(MN)$.
  \end{center}
  विस्तारित संबंध $\equal[\ext]$ असा दर्शवतो.
\end{defn}

ढोबळपणे, पदांकडे फलने म्हणून पाहिल्यास एकाच
ताज्या मुक्त चराला फलसाधक म्हणून दिल्यावर त्यांची
उपयोजने सममूल्य असतील, तर ती पदे सममूल्य मानावीत,
असे हा नियम सांगतो. येथे एखाद्या ठरावीक स्थिर आदानावर
मिळालेली समान मूल्ये पुरेशी नाहीत; नियमातील चर
दोन्ही पदांत मुक्त नसणे आवश्यक आहे.

आता $\eta$ नियमाने नेमकी विस्तारात्मकताच मिळते,
त्यापलीकडे काही नाही, हे सिद्ध करू.

\begin{thm}
  $M \equal[\ext] N$ असते, जर आणि फक्त जर $M \equal[\eta] N$.
\end{thm}

\begin{proof}
  प्रथम $\equal[\eta]$ हा संबंध विस्तारात्मकतेच्या
  नियमानुसार बंद आहे हे सिद्ध करू. म्हणजे \ext{} नियमाने
  $\equal[\eta]$ मध्ये नवीन काही जोडले जात नाही. त्यामुळे
  $\equal[\eta]$ मध्ये $\equal[\ext]$ समाविष्ट होतो;
  म्हणून $M \equal[\ext] N$ असेल, तर $M \equal[\eta] N$.

  $\equal[\eta]$ हा संबंध \ext{} नियमानुसार बंद आहे
  हे दाखवण्यासाठी लक्षात घ्या: त्या \ext{} नियमाने $M
  \equal N$ हा निष्कर्ष काढताना $Mx
  \equal[\eta] Nx$ हा पूर्वपक्ष असतो. मग $\equal$ च्या
  एका नियमानुसार $\lambd[x][Mx]
  \equal[\eta] \lambd[x][Nx]$ मिळते. दोन्ही बाजूंना
  $\eta$ नियम लावल्यावर $M \equal[\eta] N$ मिळते.

  त्याचप्रमाणे $\eta$ नियम $\equal[\ext]$ मध्ये
  समाविष्ट आहे हे सिद्ध करू. $x \notin
  FV(M)$ असलेली कोणतीही $\lambd[x][Mx]$ आणि $M$
  घेतल्यास $(\lambd[x][Mx])x \equal[\ext] Mx$ मिळते;
  म्हणून \ext{} नियमाने
  $\lambd[x][Mx] \equal[\ext] M$ मिळते.
\end{proof}



